\section{Multilinear Algebra}
\subsection{Multilinear, Symmetric, Alternating Maps}
\begin{definition}[Multilinear]\label{dfn:8.13}
	Let $\left\{ M_i \right\}_{i=1}^\ell ,P $ be $R$-modules. A map
	\[
		\phi : \prod_{i=1}^{\ell }  M_i \to P
	\]
	is $R$-multilinear if it is $R$-linear in each factor.
\end{definition}

An inductive argument shows that every $R$-multilinear map $\phi : M_1 \times \cdots \times M_{\ell }  \to P$ factors uniquely through $M_1 \otimes_R \cdots \otimes_{R} M_{\ell } $. Moreover, it follows that the tensor product is associative; that is,
\[
	M_1 \otimes _R(M_2\otimes _RM_3)\cong (M_1\otimes _RM_2)\otimes _RM_3
.\]
This allows this notation to not be abusive. We will denote
\[
	\mathbb{T} _R^{\ell }\left(M \right)  =M^{\otimes \ell }\coloneqq \bigotimes_{i=1}^\ell M
\]
for the $\ell $-fold tensor product over a singular fixed $R$-module $M$. This is called the $\ell $-th tensor power of $M$, and every $R$-multilinear map $\phi : M^{\ell }  \to P$ factors uniquely through the $\ell $th tensor power
\[\begin{tikzcd}[row sep = 2.5em, column sep = 2.5em]
	M^{\ell } \arrow[r, "\phi "]\arrow[d, "\otimes_R"']&P\\
	\mathbb{T}_R^{\ell } (M)\arrow[ur, dashed]
\end{tikzcd}\]
By convention, $\mathbb{T} _R^1(M)=M$ and $\mathbb{T} _R^0(M)=R$. We can now investigate some specific types of $R$-multilinear maps.
\begin{definition}[Symmetric/Alternating]\label{dfn:8.14}
	Let $\phi : M^{\ell }  \to P$ be $R$-multilinear.
	\begin{itemize}
		\item $\phi $ is \emph{symmetric} if a permutation of its arguments yields the same result. That is, for all $\sigma \in S_{\ell } $,
			\[
				\phi \left( m_1, \ldots, m_\ell   \right) =\phi (m_{\sigma (1), \ldots , m_{\sigma (\ell )} } )
			.\]
		\item $\phi $ is \emph{alternating} if $\phi (m_1, \ldots, m_\ell  )=0$ whenever there exist $i\neq j$ such that $m_i=m_j$.
	\end{itemize}
\end{definition}
\begin{lemma}\label{lma:8.6}
	Let $\phi : M^{\ell }  \to P$ be $R$-multilinear. If $\phi $ is alternating, then for any $\sigma \in S_\ell $, 
	\[
		\phi (m_{\sigma (1)} , \ldots ,m_{\sigma (\ell )} )=(-1)^{\sigma } \phi (m_1, \ldots, m_\ell )
	.\]
	Here, $(-1)^{\sigma }$ is 1 if $\sigma $ is an even permutation, and $-1$ otherwise. The converse holds if $2=1_R+1_R$ is a unit in $R$.
\end{lemma}
\begin{proof}
	For the first direction, it suffices to show adding a single transposition flips the sign. Thus, it suffices to show swapping the order in the case $\ell =2$ induces a sign change. By definition \ref{dfn:8.14}, $\phi (m,m)=0$ for all $m\in M$. By bilinearity, we have
	\begin{align*}
		0&=\phi (m_1+m_2,m_1+m_2)\\
		 &=\phi (m_1+m_2,m_1)+\phi (m_1+m_2,m_2)\\
		 &=\phi (m_1,m_1)+\phi (m_2,m_1)+\phi (m_1,m_2)+\phi (m_2,m_2)\\
		 &=\phi (m_2,m_1)+\phi (m_1,m_2)
	.\end{align*}
	Therefore, we are done.

	By the same logic, we can reduce the second statement to $\ell =2$. If $m_1=m_2=m$, then $\phi (m,m)=-\phi (m,m)$ implies $2\phi (m,m)=0$. Since $2$ is a unit, $\phi (m,m)=0$.
\end{proof}

\subsection{Symmetric and Exterior Powers}
The symmetric and alternating conditions come with companion universal objects, similar to those of the tensor product.
\begin{definition}[Symmetric/Exterior Powers]\label{dfn:8.15}
	The $n$th \emph{symmetric power} $\mathbb{S}_R^n(M)$ of an $R$-module $M$ is the $R$-module such that all $R$-multilinear symmetric maps $\phi : M^{n}  \to P$ factor uniquely through $\mathbb{S} _R^n(M)$:
	\[\begin{tikzcd}[row sep = 2.5em, column sep = 2.5em]
		M^n\arrow[r]\arrow[d, "\odot"']&P\\
		\mathbb{S}_R^n(M)\arrow[ur,dashed]
	\end{tikzcd}\]
	The $n$th \emph{exterior power} or \emph{wedge power} $\boldsymbol{\Lambda} _R^n\left(M \right) $ is the $R$-module such that all $R$-multilinear alternating maps $\phi : M^n \to P$ factor uniquely through $\boldsymbol{\Lambda} _R^n\left(M \right) $:
	\[\begin{tikzcd}[row sep = 2.5em, column sep = 2.5em]
		M^n\arrow[r]\arrow[d, "\wedge"']&P\\
		\boldsymbol{\Lambda} _R^n(M)\arrow[ur, dashed]
	\end{tikzcd}\]
\end{definition}
These modules come with the symmetric product $\odot$ and the exterior product, also called the wedge product, $\wedge$. To verify that such $R$-modules exist, we simply mod out by the respective relation. Let $W$ be the module generated by the set
\[
	\left\{ \bigotimes_{i=1}^n m_i \mid (m_i)\in M^{n} \land \exists i\neq j \text{ s.t. } m_i=m_j \right\} 
.\]
That is, the set of pure tensors such that $m_i=m_j$ for some $i\neq j$. We can then realize $\boldsymbol{\Lambda }_R^n(M)$ as
\[
	\boldsymbol{\Lambda}_R^n(M)\cong \frac{\mathbb{T} _R^n(M)}{W} 
.\]
The wedge product $M\wedge M$ can then be realized as the composition of maps
\[\begin{tikzcd}[row sep = 2.5em, column sep = 2.5em]
	\displaystyle{\bigwedge_{i=1}^n M}:\displaystyle{\prod_{i=1}^{n} M}\arrow[r, "\otimes "]&\displaystyle{\bigotimes_{i=1}^n M=\mathbb{T}_{R}^n(M)}\arrow[r, "\pi "]&\displaystyle{\frac{\mathbb{T} _R^n(M)}{W} }.
\end{tikzcd}\]

Similarly, we can construct $\mathbb{S} _R^n(M)$ by modding out by the set $K$ generated by tensors of the form $m_1 \otimes \ldots \otimes m_n-m_{\sigma (1)} \otimes \ldots \otimes m_{\sigma (n)} $. This identifies all pure tensors with all permutations of their indicies in the quotient. Since this sets all pure tensors equal to permutations of their indicies, any symmetric $R$-bilinear map factors uniquely through this $R$-module. The wedge product $m\wedge n$ then equals $m\otimes n$ if $m\neq n$, and $0$ otherwise. The symmetric product $m \odot n$ is then the tensor product $m\otimes n$, identified with $n\otimes m$. It can be seen as
\[\begin{tikzcd}[row sep = 2.5em, column sep = 2.5em]
	\displaystyle{\bigodot_{i=1}^n M: \prod_{i=1}^{n} M}\arrow[r, "\otimes "]&\displaystyle{\bigotimes_{i=1}^n M}\arrow[r, "\pi "]&\displaystyle{\frac{\mathbb{T} _R^n(M)}{K} }.
\end{tikzcd}\]
The images of pure tensors $m_1 \otimes \ldots \otimes m_n$ are denoted
\[
	m_1 \wedge \cdots \wedge m_n
\]
and
\[
	m_1 \odot \cdots \odot m_n
.\]

The module $\boldsymbol{\Lambda} _R^n(M)$ is generated by pure tensors of the form
\[
	m_1\wedge\cdots\wedge m_n=-m_1\wedge\cdots\wedge m_n
.\]
If $M$ is generated by a set $\left\{ e_1, \ldots, e_n  \right\} $, then $\boldsymbol{\Lambda} _R^n(M)$ is generated by all elements
\[
	e_{i_1} \wedge \cdots\wedge e_{i_\ell } 
.\]
Here, we take $1\leq i_1\leq \cdots\leq i_{\ell } \leq r$. It follows that if $M$ is free with rank $r$, then
\[
	\rank \boldsymbol{\Lambda} _R^n\left( M \right) =\binom{r}{n} 
.\]

\subsection{Very Small Detour: Graded Algebra}
\begin{definition}[Graded Ring]\label{dfn:8.16}
	A graded ring $S$ is a ring endowed with a decomposition
	\[
		S=\bigoplus_{i\geq 0}S_i
	\]
	into abelian groups $S_i$, satisfying $S_i\cdot S_j\subseteq S_{i+j} $, where $\cdot $ is the usual ring multiplication.
\end{definition}

Nonzero elements of $S_i$ are called \emph{homogeneous} of degree $i$, and it follow that the product of homogeneous elements of degree $i$ and $j$ is homogeneous of degree $i+j$. An example of a graded ring is any polynomial ring, as
\[
	R[x_1, \ldots, x_n ]=R\oplus \left\langle x_1, \ldots, x_n  \right\rangle \oplus \left\langle x^2_1, \ldots, x^2_n  \right\rangle \oplus \cdots 
.\]
\begin{definition}[Graded Module/Algebra]\label{dfn:8.17}
	Let $R$ be a graded ring. A graded $R$-module $M$ is an $R$-module $M$ such that
	\[
		M=\bigoplus_{i\geq 0}M_i
	\]
	is a decomposition of abelian groups, satisfying
	\[
		R_i\cdot M_j\subseteq M_{i+j} 
	.\]
	A graded $R$-algebra $S$ is a graded ring $S$ which is also an $R$-algebra, such that the $R$-module structure of $S$ is also graded with respect to the grading of the ring.
\end{definition}

An important example in algebraic geometry is the \emph{Rees algebra} of $I$, denoted $\operatorname{Rees}_R(I)$, and is defined as follows. Let $I\subseteq R$ be an ideal of a commutative ring. Then
\[
	\operatorname{Rees}_R(I)\coloneqq \bigoplus_{i\geq0}I^i
,\]
where $I^0=R$.

There is substantial information encoded in the grading on a ring, and as such there is significant theory behind graded algebras and modules. We will unfortunately discuss graded rings only briefly.

Graded rings form a category under ring homomorphisms, but a more interesting category is formed by only considering \emph{graded} homomorphisms; that is, homomorphisms $\phi : S \to T$ such that $\phi (S_i)\subseteq T_i$, with notation as evident. This allows for a new class of ideals to be defined, which are the kernels of graded homomorphisms.
\begin{definition}[Homogeneous Ideal]\label{dfn:8.18}
	An ideal $I$ of a graded ring $S=\bigoplus_{i\geq 0}S_i$ is homogeneous if
	\[
		I=\bigoplus_{i\geq 0}I \cap S_i
	.\]
\end{definition}
\begin{lemma}\label{lma:8.7}
	Let $S$ be a graded ring. Then the following are equivalent.
	\begin{itemize}
		\item $I$ is homogeneous;
		\item If $s\in S$ and $s=\sum_{i} s_i$ where $s_i\in S_i$ is \emph{homogeneous} for all $i$, then $s\in I$ if and only if $s_i\in I$ for all $i$;
		\item $I$ admits a generating set consisting of homogeneous elements;
		\item $I$ is the kernel of a graded homomorphism.
	\end{itemize}
\end{lemma}
\begin{proof}
	$1 \iff 2$ is immediate from definition, and $2\iff 3$ is done in the exercises.

	$2\implies 4$: Define a grading on $S/I$ by letting $(S/I)_i$ consist of cosets of elements of degree $i$ in $S$. By part 2, this is well-defined, and obviously induces a graded structure.

	$4\implies 2$. Let $\phi : S \to T$ be a graded homomorphism. Let $s=\sum_{i} s_i\in \ker \phi $, with $s_i\in S_i$ a homogeneous decomposition. Then $\sum_{i} \phi (s_i)=0$, and thus $\phi (s_i)=0$ for all $i$. Thus, $s_i\in \ker \phi $.
\end{proof}

\subsection{Tensor Algebras}
\begin{definition}[Tensor/Symmetric/Exterior Algebra]\label{dfn:8.19}
	Let $M$ be an $R$-module. The tensor algebra $\mathbb{T}_R^*(M)$ on $M$ is defined as the graded ring
	\[
		\mathbb{T}_R^*(M)\coloneqq \bigoplus_{\ell \geq 0}\mathbb{T}_R^{\ell }(M)
	.\]
	Similarly, the symmetric $\mathbb{S}_R^*(M)$ and exterior algebras $\boldsymbol{\Lambda}_R^*(M)$ are defined as
	\[
		\mathbb{S}_R^*(M)\coloneqq \bigoplus_{\ell \geq 0}\mathbb{S}_R^{\ell } (M)
	;\]
	\[
		\boldsymbol{\Lambda}_R^*(M)\coloneqq \bigoplus_{i\geq 0}\boldsymbol{\Lambda}_R^\ell (M)
	.\]
\end{definition}

These algebras are indeed $R$-algebras, as multiplication in $\mathbb{T}_R^*(M)$ can be defined as
\[
	\left( m_1 \otimes \cdots  \otimes m_k\right) \left( n_1 \otimes \cdots \otimes n_\ell  \right) = m_1 \otimes \cdots  \otimes m_k \otimes  n_1 \otimes \cdots \otimes n_\ell
.\]
Multiplication in the symmetric and exterior algebras follows by their explicit construction. The tensor algebra is not in general commutative, whereas the symmetric algebra \emph{is} commutative, and the exterior algebra is shown in the exercises to be ``skew-commutative''; that is, for $a\in \boldsymbol{\Lambda}_R^i(M)$ and $b\in \boldsymbol{\Lambda}_R^j(M)$,
\[
	a\wedge b=(-1)^{ij} b\wedge a
.\]

One important thing to show is that the quotient maps $\mathbb{T}_R^*(M)\twoheadrightarrow \mathbb{S}_R^*(M)$ and $\mathbb{T} _R^*(M)\twoheadrightarrow \boldsymbol{\Lambda}_R^*(M)$ are in fact graded ring homomorphisms; that is, their kernel is a homogeneous ideal.

\begin{lemma}\label{lma:8.8}
	Let $I_\mathbb{S} ,I_{\boldsymbol{\Lambda}} $ be the ideals generated by elements of the form $m\otimes n-n\otimes m$ and by elements of the form $m\otimes m$, respectively. Then
	\[
		\mathbb{S}_R^*(M)\cong \frac{\mathbb{T} _R^*(M)}{I_\mathbb{S} } ,\qquad \boldsymbol{\Lambda}_R^*(M)\cong \frac{\mathbb{T}_R^*(M)}{I_{\boldsymbol{\Lambda}} } 
	\]
	as graded $R$-algebras.
\end{lemma}
\begin{proof}
	This follows immediately from the explicit description of these rings, together with lemma \ref{lma:8.7}.
\end{proof}

We can now establish universal properties satisfied by these constructions.
\begin{proposition}\label{prp:8.2}
	Let $R$ be a commutative ring, and $M$ an $R$-module. Then for every $R$-algebra $T$ and every $R$-module homomorphism $\phi : M \to T$, there exists a unique $R$-algebra homomorphism $\overline{\phi } : \mathbb{T}^*_R(M) \to T$ such that the diagram
	\[\begin{tikzcd}[row sep = 2.5em, column sep = 2.5em]
		M\arrow[r, "\phi "]\arrow[d]&T\\
		\mathbb{T}_R^*(M)\arrow[ur, dashed, " \overline{\phi } "']
	\end{tikzcd}\]
	commutes.
\end{proposition}
\begin{proof}
	We have already seen that $\phi $ factors uniquely through $\mathbb{T}_R^1(M)=M$, and thus there is precisely one homomorphism $\psi : \mathbb{T}_R^*(M) \to T$ extending $ \overline{\phi } $, since they must agree on the elements of $M$, which generate $\mathbb{T} _R^*(M)$. We can simply define $\psi $ inductively such that for any $m,n\in M$, $\psi (m\otimes n)=\phi (m)\phi (n)$, and for any tensors $a,b\in\mathbb{T}_R^*(M)$, $\psi (a+b)=\psi (a)+\psi (b)$. Since $R$-linearity follows immediately, it suffices to show $\psi $ is well-defined. Indeed, this is immediate since $\phi $ is already well-defined.
\end{proof}
\begin{proposition}\label{prp:8.3}
	Let $R$ be a commutative ring, and let $M$ be an $R$-module. Then for every \emph{commutative} $R$-algebra $S$ and every $R$-module homomorphism $\psi : M \to S$, there exists a unique $R$-algebra homomorphism $\overline{\phi } : \mathbb{S}_R^*(M) \to S$ such that
	\[\begin{tikzcd}[row sep = 2.5em, column sep = 2.5em]
		M\arrow[r, "\psi "]\arrow[d]&S\\
		\mathbb{S}_R^*(M)\arrow[ur, dashed, " \overline{\psi } "']
	\end{tikzcd}\]
	commutes.
\end{proposition}
\begin{proof}
	By proposition \ref{prp:8.2}, there exists a unique $R$-algebra homomorphism $\overline{\phi } :M\to \mathbb{T}_R^*(M)$. Since $\mathbb{S}_R^*(M)$ is a quotient of the tensor algebra as a graded $R$-algebra by lemma \ref{lma:8.8}, by the universal property of quotients, $\overline{\phi } $ factors uniquely through $\mathbb{S}_R^*(M)$. Since $\mathbb{S}_R^*(M)$ is manifestly commutative by explicit description of the ideal $I_\mathbb{S} $ in lemma \ref{lma:8.8}, we are done.
\end{proof}
\begin{proposition}\label{prp:8.4}
	Let $R$ be a commutative ring, and let $M$ be and $R$-module. Then for every $R$
-algebra $A$ and every $R$-module homomorphism $\lambda : M \to A$ such that $\lambda(m)^2=0$ for all $m\in M$, there exists a unique $R$-algebra homomorphism $\overline{\lambda } : \boldsymbol{\Lambda}_R^*(M) \to A$ such that
\[\begin{tikzcd}[row sep = 2.5em, column sep = 2.5em]
	M\arrow[r, "\lambda "]\arrow[d]&A\\
	\boldsymbol{\Lambda}_R^*(M)\arrow[ur, dashed, " \overline{\lambda } "']
\end{tikzcd}\]
commutes.
\end{proposition}
\begin{proof}
	By the same argument as in proposition \ref{prp:8.3}, but applied to $\boldsymbol{\Lambda} $ instead of $\mathbb{S} $, we are done.
\end{proof}

It's easy to see the simple fact that $\mathbb{S}_R^*(R^{\oplus n} )\cong R[x_1, \ldots, x_n ]$. We can also see that these constructions $\mathbb{T} ,\mathbb{S} ,\boldsymbol{\Lambda} $ are all clearly functorial $\Rmod{R} \to \Rmod{R} $.
